\documentclass[10pt,a4paper]{amsart}
\usepackage[margin=19mm]{geometry}
\usepackage{amsmath,amssymb,mathtools}
\usepackage[scr=rsfs,cal=euler]{mathalfa}
\usepackage{xfrac}
\usepackage[unicode,hidelinks,bookmarks=false]{hyperref}
\newcommand{\ie}{i.e.\ }
\newcommand{\cf}{cf.\ }
\newcommand{\resp}{respectively\ }
\newcommand{\Subsection}{\subsection}
\providecommand{\nobreakdash}{\nobreak\hbox{-}}
\setlength{\emergencystretch}{2em}
\multlinegap0pt
% ---------------------------------------------------------------------------
% Editorial AMS wrapper prelude. The theorem environments and the mathematical macros below are
% transcribed from the paper's own preamble (cached-originals/source0.tex lines 19--52); the AMS amsart
% class, the named format packages of the pinned TeX Live runtime and this editorial formatting are not
% part of the author's text. No mathematics is changed.
% ---------------------------------------------------------------------------
\newtheorem{theorem}{Theorem}[section]
\newtheorem{prop}[theorem]{Proposition}
\newtheorem{corollary}[theorem]{Corollary}
\newtheorem{lemma}[theorem]{Lemma}
\newtheorem{conjecture}[theorem]{Conjecture}
\newtheorem{remark}[theorem]{Remark}

\theoremstyle{definition}
\newtheorem{definition}[theorem]{Definition}

\theoremstyle{example}
\newtheorem{example}[theorem]{Example}


\newcommand{\widesim}[2][1.5]{
  \mathrel{\overset{#2}{\scalebox{#1}[1]{$\sim$}}}
}

\newcommand{\brian}[1]{\color{blue}#1}
\newcommand{\james}[1]{\color{magenta}#1}

\newcommand{\boundellipse}[3]% center, xdim, ydim
{(#1) ellipse (#2 and #3)
}


\DeclareMathOperator{\mex}{mex}
\DeclareMathOperator{\myfo}{\ \widesim[2]{f_o}\ }
\DeclareMathOperator{\myfe}{\ \widesim[2]{f_e}\ }
\DeclareMathOperator{\mygo}{\ \widesim[2]{g_o}\ }
\DeclareMathOperator{\myge}{\ \widesim[2]{g_e}\ }

\allowdisplaybreaks

\begin{document}
\title{Elementary Proofs and Generalizations of Recent Congruences of Thejitha and Fathima}
\author{James A. Sellers}
\address{Department of Mathematics and Statistics, University of Minnesota Duluth, Duluth, MN 55812, USA}
\email{jsellers@d.umn.edu}
\subjclass[2020]{Primary 05A17; Secondary 11P83.}
\keywords{partitions, congruences, generating functions}
\date{Source: arXiv:2510.01572v1; distributed under CC BY 4.0}
\maketitle
\noindent\emph{Source, version and licence.} \emph{Elementary Proofs and Generalizations of Recent Congruences of Thejitha and Fathima}, by James A. Sellers. The source of this editable unit is the e-print of \textbf{version v1} of arXiv:2510.01572, https://arxiv.org/abs/2510.01572v1, whose material is distributed under the Creative Commons Attribution 4.0 International licence, http://creativecommons.org/licenses/by/4.0/, as recorded in the arXiv licence metadata for this paper and shown on that abstract page. The whole of the body below is version v1, the newest version of the paper. Full rights notice: licenses/arxiv-2510.01572-CC-BY-4.0.txt.\par\smallskip
\noindent\textit{Editorial scope.} Counted unit: original Theorem 1.3 of the article, the statement carried by the source environment \texttt{theorem} with source label \texttt{thm:TF\_mod3}, cached-originals/source0.tex lines 134--138, which reads: for all $\alpha\geq 0$ and all $n\geq 0$, $a_{5}\left(3^{2\alpha + 3}n + \frac{153\cdot 3^{2\alpha} - 1}{8}\right) \equiv 0 \pmod{3}$. It is delivered together with the author\'s own complete proof as the single primary proof body, source0.tex lines 260--343 (84 lines): the \texttt{proof} environment whose optional argument names it "Proof of Theorem \textbackslash ref\{thm:TF\_mod3\}", which carries the generating-function computation through the mod 3 dissections, derives the congruence $a_5(27n+10)\equiv a_5(3n+1)\pmod 3$ and its iterated form $a_5(81n+10)\equiv a_5(9n+1)\pmod 3$, and closes with the remark that this is the congruence the original induction uses. No proof marker is added and no mathematics is written, repaired or re-derived; the author\'s notation and the printed slips of the source are kept literal. Necessary complete context is retained uncounted, in source order: source0.tex lines 113--138, the introduction block that states the two Thejitha--Fathima theorems (Theorem 1.2 and the counted Theorem 1.3) to which this proof belongs, together with the five mod 7 congruence families quoted there; lines 145--205, the whole ``Necessary Tools\'\' section with Lemmas 2.1, 2.2 and 2.3 and their complete author proofs, including the dissection identities and the $(q;q)_\infty$ reduction that the counted proof invokes. The article\'s own numbering is reproduced by explicit counter settings: the counter \texttt{theorem}, declared by \texttt{\textbackslash newtheorem\{theorem\}\{Theorem\}[section]} and shared by prop, corollary, lemma, conjecture, remark, definition and example, is set before each retained passage, so the counted statement prints as Theorem 1.3, the introduction statements print as Theorem 1.1 and Theorem 1.2 and the three tools print as Lemmas 2.1, 2.2 and 2.3, exactly as in the article. Every \texttt{\textbackslash ref} and \texttt{\textbackslash eqref} inside every retained passage resolves inside this delivered file, and the three citations that occur in the retained passages \texttt{TF}, \texttt{Hir} and \texttt{HS2005} are delivered as the author\'s own bibliography entries transcribed verbatim from the article\'s own \texttt{thebibliography}, keeping the author\'s own printed numbers. The retained passages contain no figure environment and no graphics-inclusion command, so no artwork is redistributed. The source\'s own class is AMS \texttt{amsart} (\texttt{\textbackslash documentclass[10pt, reqno]\{amsart\}}); the e-print ships no class or style file, so no publisher class is shipped, used or redistributed, and the wrapper is the same AMS \texttt{amsart} class with the paper\'s own macros and theorem declarations transcribed from its preamble. Every declared body of this unit is an exact, unmodified span of source0.tex: no body is a bounded passage comparison, \texttt{omit} was not used anywhere, and consequently no fidelity receipt is asserted in delivery-evidence.json. The complete concatenated original source is bundled as sources/186218/source0.tex. Omitted from this editable unit and therefore absent from it are source0.tex lines 1--102; 107--112; 139--144; 206--259; 344--1153; 1157--1159; 1163--1168; 1173--1208, that is the article\'s own preamble and title block, the background discussion of the earlier literature, the sections on the elementary proofs of Theorem 1.2 and of the mod 3 divisibilities, and the whole longer second half on generalizations with their corollaries, all of which remain present in the bundled source but are not delivered here. The AMS \texttt{amsart} wrapper, the named format packages of the pinned TeX Live runtime and this editorial source, licence and scope paragraph are editorial formatting only.\par\smallskip
\setcounter{section}{1}
\setcounter{equation}{1}
\setcounter{theorem}{0}
\begin{equation}
\label{genfn_general}
\sum_{n=0}^\infty a_k(n)q^n =\frac{1}{(q^2;q^2)_\infty (q;q^2)_\infty^k} = \frac{f_2^{k-1}}{f_1^k}
\end{equation}

\setcounter{section}{1}
\setcounter{equation}{2}
\setcounter{theorem}{0}
\begin{theorem}
\label{thm:infinite_families_mod7}
For all $j\geq 0$ and all $n\geq 0$, 
\begin{align*}
a_{7j+1}(7n+5) & \equiv 0 \pmod{7},  \\
a_{7j+3}(7n+2) & \equiv 0 \pmod{7}, \\
a_{7j+4}(7n+4) & \equiv 0 \pmod{7},  \\
a_{7j+5}(7n+6) & \equiv 0 \pmod{7}, \text{\ \ and}\\
a_{7j+7}(7n+3) & \equiv 0 \pmod{7}.  
\end{align*}
\end{theorem}

Motivated by these results, Thejitha and Fathima \cite{TF} recently proved the following two sets of congruences satisfied by the function $a_5(n)$.  


\begin{theorem}
\label{thm:TF_mod5} 
For all $n\geq 0$, 
$$a_5(5n+3) \equiv 0 \pmod{5}.$$
\end{theorem}


\setcounter{section}{1}
\setcounter{equation}{2}
\setcounter{theorem}{2}
\begin{theorem}
\label{thm:TF_mod3} 
For all $\alpha \geq 0$ and all $n\geq 0$, 
$$a_{5}\left(3^{2\alpha + 3}n + \frac{153\cdot 3^{2\alpha} - 1}{8}\right)  \equiv 0 \pmod{3}.$$
\end{theorem}

\setcounter{section}{1}
\setcounter{equation}{2}
\setcounter{theorem}{3}
\section{Necessary Tools}
\label{sec:tools}

Much of our work below on the family of congruences modulo 3 relies on results found in the paper of Hirschhorn and the author \cite{HS2005}.  Using the notation of \cite{HS2005}, let 
$$D(q) = \sum_{n=-\infty}^\infty (-1)^nq^{n^2}$$ 
and 
$$Y(q) = \sum_{n=-\infty}^\infty (-1)^nq^{3n^2-2n}.$$ 
Thanks to Jacobi's Triple Product Identity \cite[(1.1.1)]{Hir}, we know that 
\begin{equation}
\label{D_prod}
D(q) = \frac{f_1^2}{f_2} 
\end{equation}
and 
\begin{equation}
\label{Y_prod}
Y(q) = \frac{f_1f_6^2}{f_2f_3}.
\end{equation}
With these in hand, we can now state the 3--dissection results that we will require in order to prove Theorem \ref{thm:TF_mod3}.  

\begin{lemma}
\label{lem:f2overf1squared_3diss} 
We have 
$$
\frac{f_2}{f_1^2} \equiv  \frac{D(q^9)^2 +2qD(q^9)Y(q^3)+q^2Y(q^3)^2}{D(q^3)} \pmod{3}.
$$
\end{lemma}
\begin{proof}
See Hirschhorn and the author \cite{HS2005}.  
\end{proof}

\begin{lemma}
\label{lem:psi_3diss}
We have 
$$
\frac{f_2^2}{f_1} = \frac{f_6f_9^2}{f_3f_{18}} +q\frac{f_{18}^2}{f_9}.  
$$
\end{lemma}
\begin{proof}
See Hirschhorn \cite[(14.3.3)]{Hir}.
\end{proof}

%\begin{lemma}
%\label{lem:f1f2_3diss}
%We have 
%$$
%f_1f_2 = \frac{f_6f_9^4}{f_3f_{18}^2} -qf_9f_{18} -2q^2\frac{f_3f_{18}^4}{f_6f_9^2}.  
%$$
%\end{lemma}
%\begin{proof}
%See Hirschhorn and the author \cite{HS2014}.
%\end{proof}


To close this section, we note a pivotal congruence result which follows from the Binomial Theorem and congruence properties of certain binomial coefficients.  
\begin{lemma} 
\label{lem:FD} 
For a prime $p$ and positive integers $a$ and $b$, we have 
$$
f_a^{bp} \equiv f_{ap}^b \pmod{p}.  
$$
\end{lemma}  

\setcounter{section}{4}
\setcounter{equation}{0}
\setcounter{theorem}{0}
\begin{proof}[Proof of Theorem  \ref{thm:TF_mod3}]
Thanks to \eqref{genfn_general} we know 
\begin{align*}
\sum_{n=0}^\infty a_5(n)q^n 
&= 
\frac{f_2^4}{f_1^5} = 
\frac{f_2^3}{f_1^3}\cdot \frac{f_2}{f_1^2} \\
&\equiv 
\frac{f_6}{f_3}\cdot \frac{f_2}{f_1^2} \pmod{3}  \\
&\equiv 
\frac{f_6}{f_3}\left( \frac{D(q^9)^2 +2qD(q^9)Y(q^3)+q^2Y(q^3)^2}{D(q^3)}  \right) \pmod{3} 
\end{align*}
using Lemma \ref{lem:f2overf1squared_3diss}.  
Thus, 
\begin{align*}
\sum_{n=0}^\infty a_5(3n+1)q^{3n+1} 
&\equiv 
2q\frac{f_6}{f_3}\left( \frac{D(q^9)Y(q^3)}{D(q^3)}  \right) \pmod{3}  \\
&= 
2q\frac{f_6f_9f_{18}}{f_3^2}
\end{align*}
using \eqref{D_prod} and \eqref{Y_prod} and simplifying the result.  This means that
\begin{align}
\sum_{n=0}^\infty a_5(3n+1)q^{n} 
&\equiv 
2f_3f_{6} \frac{f_2}{f_1^2} \pmod{3} \notag \\
&\equiv
2f_1f_2f_6 \pmod{3}.  \label{3n1_mod3}
\end{align}
Again using Lemma \ref{lem:f2overf1squared_3diss}, we see that 
\begin{align*}
\sum_{n=0}^\infty a_5(9n+1)q^{3n} 
&\equiv 
2f_3f_{6} \frac{D(q^9)^2}{D(q^3)} \pmod{3}
\end{align*}
or 
\begin{align*}
\sum_{n=0}^\infty a_5(9n+1)q^{n} 
&\equiv 
2f_1f_{2}\frac{D(q^3)^2}{D(q)} \pmod{3} \\
&= 
2f_1f_2\left( \frac{f_3^2}{f_6}\right)^2 \left( \frac{f_2}{f_1^2} \right)  = 
2\frac{f_3^4}{f_6^2}\left( \frac{f_2^2}{f_1} \right).
\end{align*}
Thanks to Lemma \ref{lem:psi_3diss}, we then know that 
\begin{align*}
\sum_{n=0}^\infty a_5(9n+1)q^{n} 
&\equiv 
2\frac{f_3^4}{f_6^2}\left( \frac{f_6f_9^2}{f_3f_{18}} +q\frac{f_{18}^2}{f_9} \right) \pmod{3}.
\end{align*}
Note that, when the above expression is expanded, it is not possible to obtain any terms involving powers of the form $q^{3n+2}$ for any $n$.  This means that, for all $n\geq 0$, 
\begin{equation*}
\label{27n19_mod3} 
a_5(9(3n+2)+1) = a_5(27n+19) \equiv 0 \pmod{3}. 
\end{equation*}
This is the $\alpha=0$ case of Theorem \ref{thm:TF_mod3}.  

We next wish to show that, for all $n\geq 0$, $a_5(81n+10)\equiv a_5(9n+1) \pmod{3}.$  Returning to the work we completed above, and noting that $a_5(9(3n+1)+1) = a_5(27n+10)$, we have 
\begin{align*}
\sum_{n=0}^\infty a_5(27n+10)q^{3n+1} 
&\equiv 
2\frac{f_3^4}{f_6^2}\left(q\frac{f_{18}^2}{f_9} \right) \pmod{3} \\
&\equiv 
2q\frac{f_3^4f_6^6}{f_3^3f_6^2}\pmod{3} \\
&= 
2qf_3f_6^4.
\end{align*}
Therefore, we know 
\begin{align*}
\sum_{n=0}^\infty a_5(27n+10)q^{n} 
&\equiv 
2f_1f_2^4 \pmod{3} \\
&\equiv 
2f_1f_2f_6 \pmod{3}.
\end{align*}
Thanks to \eqref{3n1_mod3}, we then see that, for all $n\geq 0$, 
$$
a_5(27n+10) \equiv a_5(3n+1) \pmod{3}.
$$
Replacing $n$ by $3n$ on both sides of this congruence yields $$
a_5(81n+10) \equiv a_5(9n+1) \pmod{3}
$$
and this is the congruence used by Thejitha and Fathima \cite{TF} to finalize their induction proof of Theorem \ref{thm:TF_mod3} for all $\alpha\geq 0$.  
\end{proof}

\begin{thebibliography}{99}
\bibitem[3]{Hir} M. D. Hirschhorn. \emph{The power of $q$. A personal journey}. Developments in Mathematics, 49. Springer, 2017. 
\bibitem[5]{HS2005} M. D. Hirschhorn and J. A. Sellers. An infinite family of overpartition congruences modulo 12. \emph{INTEGERS} {\bf 5} (2005), A20. 
\bibitem[8]{TF} M. P. Thejitha and S. N. Fathima. Arithmetic properties of partitions with 1-colored even parts and $r$-colored odd parts. \url{https://arxiv.org/abs/2509.24324}. \end{thebibliography} \end{document} %%%%%%%%%%%%%%%%%%%%%%%%%%%%% %%%%%%%%%%%%%%%%%%%%%%%%%%%%% %%%%%%%%%%%%%%%%%%%%%%%%%%%%% %%%%%%%%%%%%%%%%%%%%%%%%%%%%% 
\end{thebibliography}
\end{document}
